\documentclass[a4paper,11pt]{article}
\usepackage{exam-style-341}
\examinehead{341 农业综合知识三 · 网络技术分章练习}{第08章 应用层}
\begin{document}

\unitcover{341 农业综合知识三 · 网络技术分章练习}{第08章\quad 应用层}{练习（共 100 分）}{约 60 分钟}{域名系统 DNS（域名结构、域名服务器分类、递归／迭代查询）；文件传送协议 FTP（控制连接与数据连接）；电子邮件（SMTP、POP3、IMAP、MIME）；万维网 WWW（URL、HTTP、Cookie）；动态主机配置协议 DHCP。}

\begin{knowmap}
\textbf{应用层：}为应用进程提供通信服务；采用客户／服务器（C/S）或对等（P2P）方式，协议规定报文格式与交互规则\\[2pt]
\textbf{DNS：}层次树状域名结构；根 · 顶级 · 权限 · 本地域名服务器；主机到本地服务器用递归查询，本地服务器到根／顶级／权限服务器通常用迭代查询；默认 UDP 53\\[2pt]
\textbf{FTP：}控制连接（服务器端口 21，整个会话保持）＋数据连接（主动方式服务器端口 20）；两条并行的 TCP 连接，控制信息带外传送\\[2pt]
\textbf{电子邮件：}SMTP（25，发送／推送）· POP3（110）与 IMAP（143，读取／拉取）· MIME（传送非 ASCII 数据与附件）\\[2pt]
\textbf{WWW：}URL 定位资源；HTTP（80）请求／响应报文、状态码；Cookie 弥补 HTTP 的无状态性；HTML 描述页面\\[2pt]
\textbf{DHCP：}客户 UDP 68、服务器 UDP 67；发现 $\rightarrow$ 提供 $\rightarrow$ 请求 $\rightarrow$ 确认，动态分配 IP 地址并设租用期
\end{knowmap}

\begin{keyterms}
\textbf{域名系统}：把便于记忆的域名解析为 IP 地址的分布式数据库系统，采用层次树状结构。\\
\textbf{域名服务器}：存储域名与 IP 地址映射的服务器，分为根、顶级、权限、本地四类。\\
\textbf{递归查询}：主机向本地域名服务器的查询，若本地服务器不知道答案，则由它代替主机继续查询。\\
\textbf{迭代查询}：本地域名服务器向根、顶级、权限服务器的查询，对方只告诉下一步应找哪台服务器。\\
\textbf{FTP}：文件传送协议，使用 TCP 的可靠服务，在客户与服务器之间传送文件。\\
\textbf{控制连接与数据连接}：FTP 使用两条并行的 TCP 连接，控制连接传命令（端口 21），数据连接传文件内容。\\
\textbf{SMTP}：简单邮件传送协议，用于把邮件从用户代理推送到邮件服务器、以及在邮件服务器之间传送，默认端口 25。\\
\textbf{POP3}：邮局协议第 3 版，用于用户从邮件服务器拉取邮件，默认端口 110。\\
\textbf{IMAP}：互联网报文存取协议，用户可在服务器上管理邮件（如建立文件夹），默认端口 143。\\
\textbf{MIME}：通用互联网邮件扩充，使电子邮件能够传送非 ASCII 文本、图像、音频和附件。\\
\textbf{URL}：统一资源定位符，由协议、主机名（域名）、端口、路径等部分组成，用于定位网络资源。\\
\textbf{HTTP}：超文本传送协议，采用请求／响应方式，默认端口 80，是无状态协议。\\
\textbf{Cookie}：HTTP 服务器在客户端保存的小段状态信息，用于识别用户、维持会话状态。\\
\textbf{DHCP}：动态主机配置协议，自动为主机分配 IP 地址、子网掩码、默认网关和 DNS 服务器地址。
\end{keyterms}

\begin{contrast}{递归查询 / 迭代查询（每行依次为：查询发起者·后续查询由谁进行·典型使用）}
\citem{递归查询}{主机向本地域名服务器}{本地域名服务器代替主机继续查询，最终返回结果}{主机查询本地域名服务器}
\citem{迭代查询}{本地域名服务器向根、顶级、权限服务器}{被查询的服务器只返回“下一步找谁”，由本地服务器继续查询}{本地域名服务器逐级查询根、顶级、权限服务器}
\end{contrast}

\begin{contrast}{SMTP / POP3 / IMAP（每行依次为：数据方向·默认端口·主要作用）}
\citem{SMTP}{推送（发送方 $\rightarrow$ 服务器、服务器 $\rightarrow$ 服务器）}{25/TCP}{发送邮件}
\citem{POP3}{拉取（服务器 $\rightarrow$ 用户代理）}{110/TCP}{把邮件下载到本地，通常下载后删除服务器副本}
\citem{IMAP}{拉取（服务器 $\rightarrow$ 用户代理）}{143/TCP}{在服务器上管理邮件，支持多设备同步}
\end{contrast}

\begin{contrast}{FTP 控制连接 / 数据连接（每行依次为：服务器端口·传送内容·生存期）}
\citem{控制连接}{21}{传送 FTP 命令与响应，不传送文件}{整个会话期间一直保持}
\citem{数据连接}{主动方式为 20（被动方式为服务器随机端口）}{传送文件内容}{每传送一个文件建立一次，传送完即关闭}
\end{contrast}

\qtypename{一、单项选择题}{10 题，每题 3 分，共 30 分}
\q DNS 默认使用的传输层协议和端口是（\quad）
\choicesline{A. UDP 53}{B. TCP 25}{C. UDP 67}{D. TCP 80}

\q 本地域名服务器向根域名服务器查询时，通常采用（\quad）
\choicesline{A. 递归查询}{B. 迭代查询}{C. 广播查询}{D. 反向查询}

\q FTP 控制连接使用的服务器端口是（\quad）
\choicesline{A. 20}{B. 21}{C. 25}{D. 80}

\q FTP 主动方式下，服务器用于建立数据连接的熟知端口是（\quad）
\choicesline{A. 20}{B. 21}{C. 22}{D. 23}

\q SMTP 使用的默认端口是（\quad）
\choicesline{A. 21}{B. 25}{C. 80}{D. 110}

\q POP3 使用的默认端口是（\quad）
\choicesline{A. 25}{B. 53}{C. 110}{D. 143}

\q 要使电子邮件能够传送二进制附件和中文等非 ASCII 字符，需要使用的协议是（\quad）
\choicesline{A. MIME}{B. SMTP}{C. POP3}{D. SNMP}

\q HTTP 使用的默认端口是（\quad）
\choicesline{A. 21}{B. 25}{C. 53}{D. 80}

\q 浏览器访问 Web 页面时，最常用的应用层协议是（\quad）
\choicesline{A. FTP}{B. HTTP}{C. SMTP}{D. DHCP}

\q DHCP 客户使用的 UDP 端口是（\quad）
\choicesline{A. 53}{B. 67}{C. 68}{D. 69}

\qtypename{二、多项选择题}{3 题，每题 4 分，共 12 分；每题有两个或两个以上正确项}
\q 域名服务器按层次可以分为（\quad）
\choices{A. 根域名服务器}{B. 顶级域名服务器}{C. 权限域名服务器}{D. 本地域名服务器}

\q 下列属于电子邮件相关协议的有（\quad）
\choicesline{A. SMTP}{B. POP3}{C. IMAP}{D. FTP}

\q 下列属于 HTTP 请求方法的有（\quad）
\choicesline{A. GET}{B. POST}{C. HEAD}{D. SMTP}

\qtypename{三、判断题}{6 题，每题 2 分，共 12 分；正确的画“√”，错误的画“×”并写出改正后的表述}
\q DNS 只能使用 UDP 的 53 端口工作。\quad（\quad）\quad 改正：\fixblank

\q 递归查询由本地域名服务器代替主机继续查询。\quad（\quad）\quad 改正：\fixblank

\q FTP 的控制连接与数据连接是两条相互独立的 TCP 连接。\quad（\quad）\quad 改正：\fixblank

\q SMTP 是用于接收电子邮件的协议。\quad（\quad）\quad 改正：\fixblank

\q MIME 使电子邮件能够传送非 ASCII 数据和附件。\quad（\quad）\quad 改正：\fixblank

\q HTTP 是无状态协议，Cookie 可弥补其状态管理的不足。\quad（\quad）\quad 改正：\fixblank

\qtypename{四、填空题}{6 空，每空 1 分，共 6 分}
\q DNS 默认使用的端口号是\fillblanks[1.4cm]。

\q FTP 控制连接使用的端口号是\fillblanks[1.4cm]。

\q SMTP 使用的端口号是\fillblanks[1.4cm]。

\q POP3 使用的端口号是\fillblanks[1.4cm]。

\q HTTP 使用的端口号是\fillblanks[1.4cm]。

\q 电子邮件中用于传送非 ASCII 数据的扩充协议是\fillblanks[1.8cm]。

\newpage
\qtypename{五、简答题}{2 题，每题 10 分，共 20 分}
\q 简述 DNS 的域名结构、域名服务器分类，以及域名解析中递归查询与迭代查询的区别。\anslines{5cm}

\q 简述 FTP 的基本工作原理，说明控制连接与数据连接的作用，以及主动方式与被动方式的区别。\anslines{5cm}

\qtypename{六、综合题}{1 题，共 20 分}
\q 用户 A 用计算机上的邮件客户端给用户 B 发送一封带 PDF 附件的电子邮件，B 用手机邮件客户端接收。请回答：
\begin{enumerate}[leftmargin=2em, itemsep=1pt, topsep=2pt]
  \item 按顺序写出该邮件从 A 的计算机到 B 的手机所经过的环节、使用的协议及其方向（推送或拉取），并注明所用端口（8 分）。
  \item MIME 的作用是什么？为什么附件需要经过编码才能随邮件传送（4 分）？
  \item DNS 在这一过程中起什么作用（4 分）？
  \item 若 A 改为通过浏览器登录 Web 邮箱发送邮件，与使用邮件客户端方式相比有什么不同（4 分）？
\end{enumerate}
\anslines{8cm}

\end{document}
